% 图契约：八分之一波长开路支节；几何、方向和公式与正文同一约定。
\documentclass[tikz,border=5pt]{standalone}
\usepackage{fontspec,xeCJK,amsmath,amssymb}
\IfFontExistsTF{Times New Roman}{\setmainfont{Times New Roman}}{\setmainfont{TeX Gyre Termes}}
\IfFontExistsTF{Heiti SC}{\setCJKmainfont[AutoFakeBold=1.5]{Heiti SC}}{\setCJKmainfont[AutoFakeBold=1.5]{FandolHei-Regular}}
\usetikzlibrary{arrows.meta,calc,patterns}
\definecolor{ink}{HTML}{243745}
\definecolor{blue}{HTML}{275DAB}
\definecolor{green}{HTML}{177C63}
\definecolor{red}{HTML}{BD4935}
\definecolor{orange}{HTML}{AA6C16}
\definecolor{purple}{HTML}{715393}
\begin{document}
\begin{tikzpicture}[x=1cm,y=1cm,font=\fontsize{9.5}{12}\selectfont,
v/.style={-{Latex[length=2mm,width=1.3mm]},line width=.85pt},
axis/.style={-{Latex[length=1.7mm,width=1.1mm]},ink!65,line width=.55pt},
aux/.style={ink!45,dashed,line width=.55pt},
lab/.style={fill=white,inner sep=1.8pt},
card/.style={draw=ink!25,fill=blue!3,rounded corners=2pt,text width=6.1cm,minimum height=.85cm,inner sep=6pt,align=center}]
\path[use as bounding box](-8,-3.6)rectangle(8,3.65);
\node[font=\bfseries\fontsize{10.5}{13}\selectfont,text width=15.1cm,align=center]at(0.00000,3.22000){M7-048　八分之一波长开路支节};
\begin{scope}[shift={(-4,0)}]
\draw[ink,line width=.8pt](-2.50000,0.50000)--(2.00000,0.50000);
\draw[ink,line width=.8pt](-2.50000,-0.50000)--(2.00000,-0.50000);
\draw[ink](2,0.5)circle(0.06);
\draw[ink](2,-0.5)circle(0.06);
\draw[v,blue](-2.00000,1.20000)--(0.60000,1.20000);
\node[text=blue]at(-0.80000,1.60000){$V_+,I_+$};
\node[lab]at(-0.60000,0.00000){$Z_0$};
\draw[{Latex[length=1.3mm]}-{Latex[length=1.3mm]},ink!65,line width=.5pt](-2.50000,-2.20000)--(2.00000,-2.20000);
\node[lab]at(-0.25000,-2.45000){$\ell=\lambda/8$};
\draw[v,red](0.60000,-1.20000)--(-2.00000,-1.20000);
\node[text=red]at(-0.80000,-1.60000){$V_-,\ I_-=-V_-/Z_0$};
\end{scope}
\node[card]at(4.00000,1.25000){$\beta\ell=\pi/4$};
\node[card]at(4.00000,-1.05000){$Z_{\rm in}=-iZ_0\cot(\beta\ell)=-iZ_0$};
\node[text width=15.2cm,align=center]at(0.00000,-3.20000){采用 $e^{i\omega t}$；右端开路，输入呈电容性。};
\end{tikzpicture}\end{document}
